\documentclass[10pt,a4paper]{amsart}
\usepackage[margin=19mm]{geometry}
\usepackage{amsmath,amssymb,mathtools}
\usepackage[scr=rsfs,cal=euler]{mathalfa}
\usepackage{xfrac}
\usepackage[unicode,hidelinks,bookmarks=false]{hyperref}
\newcommand{\ie}{i.e.\ }
\newcommand{\cf}{cf.\ }
\newcommand{\resp}{respectively\ }
\newcommand{\Subsection}{\subsection}
\providecommand{\nobreakdash}{\nobreak\hbox{-}}
\setlength{\emergencystretch}{2em}
\multlinegap0pt
\usepackage{bm}
\usepackage{bbm}
\usepackage{dsfont}
\usepackage{xspace}
\usepackage{cancel}
\usepackage{mathtools}
\usepackage{amsthm}
\makeatletter
\@ifundefined{theorem}{\newtheorem{theorem}{Theorem}}{}
\@ifundefined{prop}{\newtheorem{prop}[theorem]{Proposition}}{}
\makeatother
\title[Stability of Homogeneous minimal hypersurfaces in the Page s]{Stability of Homogeneous minimal hypersurfaces in the Page space and $Y^{p,q}$ Sasaki-Einstein manifolds}
\author{Natalia Gherghel, Hari K. Kunduri}
\date{Source: arXiv:2511.05447v1 (CC BY 4.0)}
\begin{document}
\maketitle

\noindent\emph{Source and licence.} Stability of Homogeneous minimal hypersurfaces in the Page space and $Y^{p,q}$ Sasaki-Einstein manifolds. Version arXiv:2511.05447v1 (CC BY 4.0), distributed under the Creative Commons Attribution 4.0 International licence, as recorded in the arXiv licence metadata for this exact version and shown on the abstract page https://arxiv.org/abs/2511.05447v1. Full rights notice: licenses/arxiv-2511.05447-CC-BY-4.0.txt.\par\smallskip

\noindent\textit{Editorial scope.} Counted unit: the Prop whose source label is \texttt{None} in the article (source0.tex lines 375--376, source label \texttt{None}), delivered together with the article's own complete original proof of it (source0.tex lines 377--392, one \texttt{proof} environment of 16 lines). No mathematics is written, repaired or re-derived: statement and proof are the article's own text line for line. Required context retained uncounted: none. Every definition, display and label of those lines is retained unchanged. Omitted from this package and not redistributed: 1--374, 393--638. Nothing inside a retained excerpt is omitted or altered: every retained body is delivered exactly as the article's lines are written, so no omission receipt of any kind accompanies this package. Citations: the retained excerpts carry no citation command at all, so no bibliography entry of the article is transcribed and no reference is redistributed. Reference adaptation: none was needed and none was made; every reference inside the delivered unit resolves inside it, so no unresolved reference or citation is produced. Printed slips kept literal: no printed word, symbol, hypothesis or number of the counted statement or of its proof was altered. Layout and class: the article's own class is \texttt{amsart} with packages including amsmath, amssymb, upgreek, tensor, xcolor, url, xy, longtable, graphicx, float, bm, enumerate, subfig, tikz; the wrapper is the lane's pinned \texttt{amsart} editorial wrapper at 10pt on A4, which restates the theorem-like environments the retained text needs and adds the article's own macro definitions that the retained text needs (none), with extra wrapper packages (bm, bbm, dsfont, xspace, cancel, mathtools). The delivered numbering is the unit's own. Assets: no retained span contains a figure environment, a graphics-inclusion command, a PDF-page inclusion, a table or a TikZ picture; no image file is copied into this package, the package declares no asset dependency and no image file is redistributed with it. Licence: version arXiv:2511.05447v1 of this article is distributed under the Creative Commons Attribution 4.0 International licence, as recorded in the arXiv licence metadata for this exact version and shown on the abstract page https://arxiv.org/abs/2511.05447v1. A full rights notice is delivered at licenses/arxiv-2511.05447-CC-BY-4.0.txt. Characters: every retained excerpt is pure ASCII, so the delivered engine, XeLaTeX with its default Latin Modern text font, reports no missing character.
\begin{prop} There is a single minimal hypersurface of the above form, corresponding to a level set $y = y_m$ where $y_m \in (y_1,0)$. 
\end{prop}

\begin{proof} The determinant of the induced metric $h_{ij}$ is given by
\begin{equation}
\det h = \frac{(1-y)^2}{324} q(y) w(y) \sin^2\theta.
\end{equation} The mean curvature of the hypersuface is then 
\begin{equation}
    H_\Sigma = \text{Tr}_h K = h^{ij}K_{ij} = \frac{1}{2\sqrt{g_{yy}}} \partial_y (\log \det h).
\end{equation} Straightforward computation shows that the condition $H_\Sigma=0$ is equivalent to solving the cubic
\begin{equation}
    H(y):=b + 6 y - 15 y^2 + 8 y^3 =0
\end{equation}  The discriminant of $H(y)$ is 
\begin{equation}
    -108 \left(16 b^2-5b - 11\right) = 108 \cdot  \left(16 b + 11\right)(1-b)
\end{equation} In the range $0 < b < 1$, this quantity is strictly positive. Thus, there are three distinct real roots $y^*_i$, $i=1,2,3$. It remains to determine which of these roots lie in the range $(y_1, y_2)$ and thus constitute a minimal surface. From Descartes' rule of signs, $H(y)$ can have two or zero positive roots. It is easy to check that the cubic $H(y)$ has a local maximum at $y = 1/4$ and a local minimum at $y = 1$. In addition $H(0) >0$, $H(1/4) > 0$, and $H(1) = b-1 < 0$. It follows that $H(y)$ must have positive roots $y_2^*, y_3^*$ at some $1/4 < y_2^* < 1$, $y_3^* > 1$ and a single negative root $y_1^*$. Clearly $y_3^*$ cannot lie inside $(y_1,y_2)$. We can also see that the cubic $Q(y)$ has a local maximum at $y = 0$ and a local minimum at $y = 1$. Thus (as expected) it will be positive between one negative root $y_1 <0$ and one positive root $y_2 \in (0,1)$, with a third positive root $y_3>1$. 

At any root $y_i^*$ of $H(y)$,  $Q(y) = b - 3y^2 + 2y^3 = -6 y_i^{*}(1 - y_i^{*})^2$. Observe that for $y_1^* < 0$, this cubic is positive; but for $y^*_2 > 0$ it will be negative. In particular, $g_{yy} < 0$ and the induced metric $h$ is not positive definite (since $q(y^*_2) < 0$). It follows that only the negative root $\bar{y}:=y_1^* \in (y_1,y_2)$ of $H(y)$ gives rise to a minimal surface. 
\end{proof}

\end{document}
